% 第 IX 讲，第10—11节；原 PDF 448—473 页。
\SourcePage{448}
\SGAsection{10}{单值性配对的性质；整性与正性定理}
本节首先给出\Term{monodromy-pairing}{单值性配对}的一些本质上平凡的\Term{form}{形式}性质；那些稍嫌繁琐的核查细节留给读者。随后，我们陈述一个关于这些\Term{pairing}{配对}的较深\Term{theorem}{定理}（10.4），并说明如何将其证明归结为 Jacobian 的一个特殊情形。这个情形将在第12节用\Term{transcendental}{超越}方法处理；第12节在逻辑上也独立于第11节。

\textbf{10.1。\Term{transport-of-structure}{结构搬运}。}由\Term{monodromy-pairing}{单值性配对} (9.1.2) 的定义可知，这个\Term{pairing}{配对}与情形 $(S,A_K,A'_K)$ 的\Term{transport-of-structure}{结构搬运}相容。

\textbf{10.2。\Term{symmetry}{对称性}。}如果从 $A_K$ 的\Term{dual-scheme}{对偶概形} $A'_K$ 出发，对它应用第9节的构造，就须引入 $A'_K$ 的\Term{dual-scheme}{对偶概形} $(A'_K)'$；由\Term{biduality}{双对偶定理}（[5] 或第 VIII 讲3.2），后者\Term{canonical}{典范}\Term{isomorphism}{同构}于 $A_K$。以 $u_\ell^{A_K}$ 表示相对于 $A_K$ 的\Term{monodromy-pairing}{单值性配对} (9.1.2)，则同样得到一个\Term{pairing}{配对}
\begin{equation}\tag{10.2.1}
 u_\ell^{A'_K}:\underline M'_\ell\otimes\underline M_\ell\longrightarrow\mathbb Z_\ell.
\end{equation}
我断言，后一个\Term{pairing}{配对}由 $u_\ell^{A_K}$ 通过交换\Term{isomorphism}{同构}
\[
 s:\underline M_\ell\otimes\underline M'_\ell\xrightarrow{\sim}\underline M'_\ell\otimes\underline M_\ell
\]
得到，即有公式
\SourcePage{449}
\begin{equation}\tag{10.2.2}
 u_\ell^{A'_K}=u_\ell^{A_K}\circ s.
\end{equation}
考虑到\Term{pairing}{配对} $\varphi_\ell$ 已知的\Term{anticommutativity}{反交换性}（[5] 或第 VIII 讲2.2.11），\Term{symmetry}{对称}公式 (10.2.2) 可用 $A_K$ 的一个给定\Term{polarization}{极化} $\xi$ 改述如下。它定义一个\Term{homomorphism}{同态} $\widetilde\xi:A_K\longrightarrow A'_K$，因而给出相应\Term{homomorphism}{同态}（$t$ 是 torique，即“\Term{torus}{环面}的”的首字母）
\begin{equation}\tag{10.2.3}
 \xi_t:\underline M\longrightarrow\underline M',
\end{equation}
从而可以构成\Term{pairing}{配对}
\begin{equation}\tag{10.2.4}
 \varphi^t_{\xi,\ell}:\underline M_\ell\otimes\underline M_\ell\longrightarrow\mathbb Z_\ell,
 \qquad \varphi^t_{\xi,\ell}(x,y)=u_\ell^{A_K}(x,\xi_t(y)).
\end{equation}
上述\Term{form}{形式} $\varphi^t_{\xi,\ell}$ 是\Term{symmetric}{对称}的：
\begin{equation}\tag{10.2.5}
 \varphi^t_{\xi,\ell}(x,y)=\varphi^t_{\xi,\ell}(y,x).
\end{equation}
事实上，对 $A_K$ 的\Term{neron-severi-group}{Néron--Severi 群}的任一元素 $\xi$（不必是\Term{polarization}{极化}），这个陈述都成立；不过，只有当 $\xi$ \Term{nondegenerate}{非退化}，即定义一个\Term{isogeny}{同源} $A_K\longrightarrow A'_K$，从而也给出一个\Term{isogeny}{同源} (10.2.3) 时，\Term{symmetry}{对称}关系 (10.2.5) 才\Term{equivalence}{等价}于\Term{symmetry}{对称}关系 (10.2.2)。而且，当 $\xi$ 是\Term{polarization}{极化}时，我们将在下面10.4 b) 中对 (10.2.5) 作出强得多的说明。

最后指出解释这个\Term{symmetry}{对称}关系的第三种方式。考虑 $K$ 上一对具有\Term{semistable-reduction}{半稳定约化}的\Term{abelian-scheme}{阿贝尔概形} $(A_K,\overline A_K)$ 上的一个\Term{biextension}{双扩张}（即\Term{divisorial-correspondence}{除子对应}）$\xi$，它定义一个\Term{homomorphism}{同态}
\[
 A_K\longrightarrow\overline A'_K.
\]
\SourcePage{450}
通过\Term{neron-model}{Néron 模型}，它便在\Term{special-fiber}{特殊纤维}的最大\Term{torus}{环面}之间定义一个\Term{homomorphism}{同态}，从而在\Term{character-group}{特征标群}之间定义一个反向的\Term{homomorphism}{同态}，于是得到
\begin{equation}\tag{10.2.6}
 \xi_t:\underline{\overline M}\longrightarrow\underline M',
\end{equation}
其中 $\underline{\overline M}$ 由 $\overline A_K$ 定义，正如 $\underline M$ 由 $A_K$ 定义一样。因此可以构成\Term{pairing}{配对}
\begin{equation}\tag{10.2.7}
 \varphi^t_\xi:\underline M_\ell\otimes\underline{\overline M}_\ell\longrightarrow\mathbb Z_\ell,
 \qquad \varphi^t_\xi(x,y)=\varphi^t(x,\xi_t(y)).
\end{equation}
顺便指出，当具有\Term{semistable-reduction}{半稳定约化}的\Term{abelian-scheme}{阿贝尔概形} $A_K$ 变化时，\Term{twisted-constant-group}{扭常值群} $\underline M$ 是关于 $A_K$ 的\Term{covariant-functor}{协变函子}；于是，在 $(K,S)$ 上，赋予具有\Term{semistable-reduction}{半稳定约化}的\Term{abelian-scheme}{阿贝尔概形}的\Term{biextension}{双扩张}以\Term{pairing}{配对} (10.2.7) 的构造，显然与\Term{biextension}{双扩张}的\Term{inverse-image}{逆像}相容。特别地，当 $\xi$ 是 $(A_K,A'_K)$ 的\Term{canonical}{典范}（“Weil”）\Term{biextension}{双扩张}时，\Term{pairing}{配对} (10.2.7) 就是\Term{monodromy-pairing}{单值性配对} (9.2.1)。现在，\Term{symmetry}{对称}关系可表述如下：考虑 $\xi$ 的交换\Term{biextension}{双扩张} $s\xi$，它是 $(\overline A_K,A_K)$ 的一个\Term{biextension}{双扩张}，因而给出\Term{pairing}{配对}
\[
 \varphi^t_{s\xi}:\underline{\overline M}_\ell\otimes\underline M_\ell\longrightarrow\mathbb Z_\ell;
\]
此时有\Term{symmetry}{对称}关系
\begin{equation}\tag{10.2.8}
 \varphi^t_{s\xi}=\varphi^t_\xi\circ s,
 \qquad s:\underline{\overline M}_\ell\otimes\underline M_\ell\xrightarrow{\sim}
 \underline M_\ell\otimes\underline{\overline M}_\ell,
\end{equation}
其中 $s$ 仍是交换\Term{isomorphism}{同构}。读者可以将此关系与第 VIII 讲2.2.11中的\Term{antisymmetry}{反对称}关系比较。

\SourcePage{451}
\textbf{10.3。\Term{trait}{离散赋值环谱}的变换。}现设
\begin{equation}\tag{10.3.1}
 g:S'\longrightarrow S
\end{equation}
为 Hensel \Term{trait}{离散赋值环谱}之间的\Term{dominant-morphism}{占优态射}，对应于\Term{discrete-valuation-ring}{离散赋值环}之间的\Term{injective}{单射}\Term{local-homomorphism}{局部同态}
\[
 V\longrightarrow V'.
\]
照常置
\begin{equation}\tag{10.3.2}
 e(S'/S)=\operatorname{long}_{V}(V'/\mathfrak m V'),
\end{equation}
其中 $\mathfrak m$ 是 $V$ 的\Term{maximal-ideal}{极大理想}。已知 (10.3.1) 在\Term{inertia-group}{惯性群}的“\Term{tame}{温和}”部分上诱导的\Term{homomorphism}{同态}
\begin{equation}\tag{10.3.3}
 g_*:I'_t\xrightarrow{\sim}\prod_{\ell\ne p}\mathbb Z_\ell(1)(k')
 \longrightarrow I_t=\prod_{\ell\ne p}\mathbb Z_\ell(1)(k)
\end{equation}
由下式给出：
\begin{equation}\tag{10.3.4}
 g_*(x)=e(S'/S)x,
\end{equation}
这里用显然的\Term{isomorphism}{同构} $\mathbb Z_\ell(1)(k)\xrightarrow{\sim}\mathbb Z_\ell(1)(k')$ 来认同 (10.3.3) 的两端（第 I 讲\SGAsourcegap{此处第 I 讲的引注编号为空白}）。

另一方面，对于在 $S$ 上具有\Term{semistable-reduction}{半稳定约化}的 $K$ 上\Term{abelian-scheme}{阿贝尔概形} $A_K$，由定义（特别是3.4）立刻可知，\Term{twisted-lattice}{扭格} $\underline M,\underline M'$ 的构造与上述任意\Term{base-change}{基变换} $S'\longrightarrow S$ 交换。因此可以将 $A_{K'}$（$K'$ 为 $V'$ 的\Term{fraction-field}{分式域}）定义的\Term{pairing}{配对} $u_\ell^{A_{K'}}$ 视为一个\Term{pairing}{配对}
\[
 u_\ell^{A_{K'}}:\underline M_{S'}\otimes\underline M'_{S'}\longrightarrow\mathbb Z_{\ell S'},
\]
\SourcePage{452}
并将它与从 $u_\ell^{A_K}$ 通过\Term{base-change}{基变换}得到的\Term{pairing}{配对} $u_{\ell S'}^{A_K}$ 比较。当 $\ell\ne p$ 时，第9.2节给出的\Term{pairing}{配对} $u_\ell$ 的构造与公式 (10.2.3) 立刻表明
\begin{equation}\tag{10.3.5}
 u_\ell^{A_{K'}}=e(S'/S)u_\ell^{A_K}.
\end{equation}
事实上，这个公式在 $\ell=p$ 时也成立。为此回到9.3至9.6的\Term{generic}{一般}构造，并归结为如下相容性陈述：对于\Term{strictly-local}{严格局部}的\Term{trait}{离散赋值环谱}的变换 $S'\longrightarrow S$，\Term{isomorphism}{同构}9.4.4 e) 满足与 (10.3.5) 类似的相容性。（注意，所讨论的\Term{isomorphism}{同构}来自由\Term{norm}{范数}给出的\Term{isomorphism}{同构} $K^*/V^*\xrightarrow{\sim}\mathbb Z$，因而给出如下\Term{commutative-diagram}{交换图}：
\[
\begin{tikzcd}
 K^*/V^*\arrow[r,"\sim"]\arrow[d]&\mathbb Z\arrow[d,"e(S'/S)"]\\
 K'^*/V'^*\arrow[r,"\sim"]&\mathbb Z.
\end{tikzcd}
\]
）

下面是关于\Term{monodromy-pairing}{单值性配对}的一个较深结果；我们将在第12节利用本讨论班以后才证明的一些结果，给出它的\Term{transcendental}{超越}证明。（在本讨论班的后续部分，我们既不会使用10.4，也不会使用本讲的其他结果。）

\SourcePage{453}
\begin{SGAresult}{\Term{theorem}{定理}}{10.4}
设 $S$ 为 Hensel \Term{trait}{离散赋值环谱}，$A_K$ 为 $K$ 上在 $S$ 上具有\Term{semistable-reduction}{半稳定约化}的\Term{abelian-scheme}{阿贝尔概形}。考虑像9.1（参见5.5.2）那样与之相伴的\Term{lattice}{格} $\underline M,\underline M'$。有如下结论：

 a) 存在唯一的\Term{pairing}{配对}
\begin{equation}\tag{10.4.1}
 u:\underline M\otimes\underline M'\longrightarrow\mathbb Z_S,
\end{equation}
使得对每个\Term{prime-number}{素数} $\ell$，$u\otimes_{\mathbb Z}\mathbb Z_\ell$ 都是 (9.1.2) 的\Term{monodromy-pairing}{单值性配对} $u_\ell$。这个\Term{pairing}{配对}是\Term{separating-pairing}{分离的}（séparant）。

 b) 设 $\xi$ 是 $A_K$ 的一个\Term{polarization}{极化}，因而定义\Term{homomorphism}{同态} (10.2.3)
\[
 \xi_t:\underline M\longrightarrow\underline M',
\]
并考虑\Term{form}{形式}
\begin{equation}\tag{10.4.2}
 \varphi^t_\xi:\underline M\otimes\underline M\longrightarrow\mathbb Z_S,
 \qquad \varphi^t_\xi(x,y)=u(x,\xi_t(y)).
\end{equation}
此\Term{form}{形式}\Term{symmetric}{对称}、\Term{positive-form}{正}且\Term{nondegenerate}{非退化}。
\end{SGAresult}
当然，用 $\underline M$ 在 $S$ 上某个给定\Term{geometric-point}{几何点}处的\Term{geometric-fiber}{几何纤维} $M_0$ 来表述，上述断言就是说，相应的\Term{form}{形式}
\[
 M_0\otimes M_0\longrightarrow\mathbb Z
\]
\Term{symmetric}{对称}、\Term{positive-form}{正}且\Term{nondegenerate}{非退化}。至于\Term{pairing}{配对} (10.4.1)，我们仍称它为\Term{monodromy-pairing}{单值性配对}。

\SourcePage{454}
\textbf{10.5。将10.4归结为一个特殊情形。}

\textbf{10.5.1。}满足所述条件的 $u$ 的唯一性显然成立，即使只要求对一个 $\ell$ 满足条件也是如此。另一方面，存在性只是说，每个 $u_\ell$ 实际上都把 $\underline M\otimes\underline M'$ 映入\SGAsourcegap{10.5.1此处的值域公式因原页磨损而缺失}，并且如此得到的\Term{homomorphism}{同态} $\underline M\otimes\underline M'\longrightarrow\mathbb Z_S$ 与 $\ell$ 无关。事实上，只须对 $u_\ell\otimes_{\mathbb Z_\ell}\mathbb Q_\ell$ 及某个 $u_{\mathbb Q}:\underline M_{\mathbb Q}\otimes\underline M'_{\mathbb Q}\longrightarrow\mathbb Q_S$ 证明相应断言就够了。

现假定，在像9.3那样作\Term{base-change}{基变换} $S'\longrightarrow S$ 后，已经证明存在性。由 (10.3.5) 立即得到一个\Term{homomorphism}{同态}
\[
 u:\underline M\otimes\underline M'\longrightarrow
 \frac{1}{e(S'/S)}\mathbb Z_S\subset\mathbb Q_S,
\]
它与各 $u_\ell\otimes_{\mathbb Z_\ell}\mathbb Q_\ell$ 相容，由此得到 a)。至于 b)，显然，若它对 $A_{K'}$ 成立，就也对 $A_K$ 成立。

因此，证明10.4时，可以任意作\Term{dominant-morphism}{占优}的\Term{trait}{离散赋值环谱}变换 $S'\longrightarrow S$。特别地，可以假定 $S$ \Term{complete}{完备}且其\Term{residue-field}{剩余域}\Term{algebraically-closed}{代数闭}。此时 $\underline M,\underline M'$ 都是\Term{constant-group-scheme}{常值群概形}，分别取值于\Term{finite-type}{有限生成}\Term{free-module}{自由} $\mathbb Z$-模 $M,M'$，于是可以方便地用后者改述 a) 和 b)。

\textbf{10.5.2。}现假定 $A_K$ \Term{isomorphism}{同构}于 $k$ 上一个具有\Term{semistable-reduction}{半稳定约化}的\Term{abelian-scheme}{阿贝尔概形} $\overline A_K$ 的\Term{direct-factor}{直因子}，且10.4已经对后者证明。考虑到\Term{lattice}{格} $M,M'$ 和\Term{monodromy-pairing}{单值性配对}在具有\Term{semistable-reduction}{半稳定约化}的\Term{abelian-scheme}{阿贝尔概形}的\Term{direct-sum}{直和}下显然的\Term{functorial}{函子性}，立刻可知10.4也对 $A_K$ 成立。
\SourcePage{455}
由10.5.1也立刻可知，10.4是否成立只取决于 $A_K$ 的\Term{isogeny}{同源}类。将此与上面的论述及“Poincaré \Term{reducible}{可约}化\Term{theorem}{定理}”[15，\SGAsourcegap{原书此处参考文献15后的定位引注为空白}] 结合，就知道：若 $A_K$ 是某个满足10.4的\Term{abelian-scheme}{阿贝尔概形} $\overline A_K$ 的\Term{quotient}{商}，则 $A_K$ 也满足10.4。

\textbf{10.5.3。}现在利用如下熟知事实 [31, th.~11 (p.~20)]：必要时对 $K$ 作一个\Term{finite}{有限}\Term{extension}{扩张}（由10.5.1，这是允许的），每个\Term{abelian-scheme}{阿贝尔概形} $A_K$ 都是 $X$ 上一条\Term{proper}{固有}、\Term{smooth}{光滑}、\Term{geometrically-connected}{几何连通}\Term{curve}{曲线} $X_K$ 的\Term{jacobian}{Jacobian}（$=\operatorname{Pic}^0_{X_K/K}$）的\Term{quotient}{商}。考虑到10.5.2，便归结为对这样的\Term{jacobian}{Jacobian}证明10.4。

\textbf{10.5.4。}由 Artin--Mumford 的一个\Term{theorem}{定理} [3]（已在第 V 讲\SGAsourcegap{此处第 V 讲的引注编号为空白}中引用并使用），必要时再作 $K$ 的一个\Term{finite}{有限}\Term{extension}{扩张} $K'$，可以假定 $X_K$ 延拓为 $S$ 上一个\Term{projective}{射影}、\Term{flat}{平坦}的\Term{regular}{正则}\Term{scheme}{概形}，且其\Term{geometric-special-fiber}{几何特殊纤维} $X_{\overline s}$ 的\Term{singular-point}{奇点}只有\Term{ordinary-quadratic-singularity}{普通二次奇点}（在这里即两条\Term{branch}{分支}的\Term{normal-crossing}{正规交叉}）。在这种情形，我们将在下面12.5中，通过用此情形的几何数据完全描述\Term{monodromy-pairing}{单值性配对} (10.4.1)，证明10.4 a)，同时在\Term{jacobian}{Jacobian}的\Term{canonical}{典范}\Term{polarization}{极化}情形核实10.4 b)。

\SourcePage{456}
\textbf{10.5.5。}还须说明：对于给定的具有\Term{semistable-reduction}{半稳定约化}的 $A_K$，若10.4 a) 成立，且10.4 b) 对 $A_K$ 的某个给定\Term{polarization}{极化} $\xi_0$ 成立，则该结论对 $A_K$ 的每个\Term{polarization}{极化} $\xi$ 都成立。事实上，设 $\widetilde\xi_0,\widetilde\xi$ 是由 $\xi_0,\xi$ 定义的\Term{homomorphism}{同态}
\[
 \widetilde\xi_0,\widetilde\xi:A_K\longrightarrow A'_K;
\]
它们都是\Term{isogeny}{同源}，且有
\[
 \widetilde\xi=\widetilde\xi_0\alpha={}^t\!\alpha^*\widetilde\xi_0,
 \qquad \alpha\in\operatorname{End}(A_K)\otimes_{\mathbb Z}\mathbb Q,
\]
其中 $\alpha\mapsto\alpha^*$ 是由\Term{polarization}{极化} $\xi_0$ 定义的\Term{algebra}{代数} $\operatorname{End}(A_K)\otimes_{\mathbb Z}\mathbb Q$ 的\Term{canonical}{典范}\Term{involution}{对合}；${}^t(\alpha^*)$ 中的上标 $t$ 表示\Term{transpose}{转置}（将 $A_K$ 的\Term{endomorphism}{自同态}变为 $A'_K$ 的\Term{endomorphism}{自同态}）。另一方面，已知 [15, Chap.~V, th.~3]
\[
 \alpha=\sum_i\beta_i^*\beta_i,
 \qquad \beta_i\in\operatorname{End}(A_K)\otimes_{\mathbb Z}\mathbb Q,
\]
从而
\[
 \widetilde\xi=\sum_i{}^t\!\beta_i\widetilde\xi_0\beta_i.
\]
或者，用所考察的\Term{polarization}{极化}所对应的\Term{biextension}{双扩张}（即\Term{divisorial-correspondence}{除子对应}）（第 VIII 讲4.14）来表述，便有
\begin{equation}\tag{10.5.5.1}
 \delta(\xi)=\sum_i\delta(\xi_0)\circ(\beta_i,\beta_i).
\end{equation}
由 (10.2.7) 后的说明，这个公式蕴含（采用 (10.4.2) 的记号）
\begin{equation}\tag{10.5.5.2}
 \varphi^t_\xi=\sum_i\varphi^t_{\xi_0}\circ(\beta_i,\beta_i)
 \in\operatorname{Hom}(M_{\mathbb Q}\otimes M_{\mathbb Q},\mathbb Q).
\end{equation}
\SourcePage{457}
这确实说明，若 $\varphi^t_{\xi_0}$ \Term{symmetric}{对称}且\Term{positive-form}{正}，则 $\varphi^t_\xi$ 也如此。

最后，$\varphi^t_{\xi_0}$ \Term{nondegenerate}{非退化}显然\Term{equivalence}{等价}于\Term{pairing}{配对} (10.4.1) 是\Term{separating-pairing}{分离的}，因为 $\xi_{0t}:\underline M\otimes_{\mathbb Z}\mathbb Q\longrightarrow\underline M'\otimes_{\mathbb Z}\mathbb Q$ 是\Term{isomorphism}{同构}，而 $\widetilde\xi_0$ 是\Term{isogeny}{同源}；因此这也与所选的\Term{polarization}{极化}无关。还可以指出，要核实 (10.4.1) 是\Term{separating-pairing}{分离的}，只须对某一个\Term{prime-number}{素数} $\ell$ 核实\Term{monodromy-pairing}{单值性配对} $u_\ell$ (9.1.2) 满足同一断言。取 $\ell\ne p$，便可利用9.2对 $u_\ell$ 的描述；此时 $u_\ell$ 在左侧是\Term{separating-pairing}{分离的}这一事实显然成立，因\Term{rank}{秩}相同，它在两侧都是\Term{separating-pairing}{分离的}。

\SGAsection{11}{Néron 模型的连通分支：对偶关系与渐近行为}
\textbf{11.0。}本节旨在研究第1节已经考察过的\Term{finite}{有限}\Term{etale}{平展} $k$-群
\begin{equation}\tag{11.0.1}
 \Phi_0=A_0/A_0^0.
\end{equation}
对于每个\Term{prime-number}{素数} $\ell$，以 $\Phi_0(\ell)$ 表示它的 $\ell$-\Term{primary-component}{准素分量}。由公式
\begin{equation}\tag{11.0.2}
 \Phi_0=\coprod_\ell\Phi_0(\ell),
\end{equation}
知道 $\Phi_0$ 就\Term{equivalence}{等价}于知道各个 $\Phi_0(\ell)$。
\SourcePage{458}
注意，群 $\Phi_0,\Phi_0(\ell)$ 分别由通常意义下的\Term{finite-group}{有限群} $\Phi_0(\overline k),\Phi_0(\ell)(\overline k)$ 以及 $\pi_0=\operatorname{Gal}(\overline k/k)$ 在它们上的自然\Term{action}{作用}决定。另一方面，$\Phi_0(\overline k)$ 显然也是 $A_{\widetilde K}$ 的\Term{neron-model}{Néron 模型}的\Term{connected-components}{连通分支群}；这里 $\widetilde K$ 是 $S$ 关于所选的 $k$ 的\Term{separable-closure}{可分闭包} $\overline k$ 的\Term{strict-henselization}{严格 Hensel 化}的\Term{fraction-field}{分式域}。因此，研究 $\Phi_0$ 时可归结为 $S$ \Term{strictly-local}{严格局部}的情形；$\pi_0\simeq\operatorname{Gal}(\widetilde K/K)$ 在 $\Phi_0(\overline k)$ 上的\Term{action}{作用}由\Term{transport-of-structure}{结构搬运}自动恢复。也可通过\Term{base-change}{基变换} $\widehat S\longrightarrow S$，进一步归结为 $S$ \Term{complete}{完备}的情形，因为\Term{neron-model}{Néron 模型}的构造与此\Term{base-change}{基变换}交换 [21, th.~3.4 b)]。

\textbf{11.1。}现假定 $S$ \Term{strictly-local}{严格局部}，于是可将 $\Phi_0$ 认同为通常意义下的\Term{finite-group}{有限群} $\Phi_0(k)$，并得到
\begin{equation}\tag{11.1.1}
 \begin{aligned}
 \Phi_0(k)&\xrightarrow{\sim}A_0(k)/A_0^0(k)
 \xrightarrow{\sim}A(S)/A^0(S)\\
 &\xrightarrow{\sim}A(K)/A(K)^0.
 \end{aligned}
\end{equation}
这里第一个等式来自 $A_0$ 在 $k$ 上\Term{smooth}{光滑}；第二个来自“Hensel \Term{lemma}{引理}”，因为 $A$ 在 $S$ 上\Term{smooth}{光滑}，而 $S$ 是\Term{henselian}{Hensel}的；最后，第三个来自关系 $A(S)\xrightarrow{\sim}A(K)$ 以及我们的约定
\begin{equation}\tag{11.1.2}
 A(K)^0\mathrel{\mathop=^{\mathrm{def}}}A^0(S).
\end{equation}
因此 $A(K)^0$ 是 $A(K)$ 的一个自然的\Term{finite-index-subgroup}{有限指数子群}（不要将它与 $A^0(K)=A(K)$ 混淆！）。这些\Term{isomorphism}{同构}
\SourcePage{459}
(11.1.1) 显然关于 $A$ 具有\Term{functorial}{函子性}，并与 $(S,A_K)$ 的\Term{transport-of-structure}{结构搬运}相容，特别地，与可能出现的\Term{galois-group}{Galois 群}的\Term{action}{作用}相容。

现设 $\ell$ 为不同于 $p$ 的\Term{prime-number}{素数}，于是 $A^0$ 是 $\ell$-\Term{divisible}{可除}的；更确切地说，$\ell\operatorname{id}_{A^0}$ 是\Term{etale}{平展}\Term{surjective}{满射}\Term{morphism}{态射}，从而在 $A^0(S)=A(K)^0$ 上诱导一个\Term{surjective}{满射}\Term{endomorphism}{自同态}。故后者是一个 $\ell$-\Term{divisible-group}{可除群}。立即可知，对每个与 $p=\operatorname{car}(k)$ 互素的整数 $n\ne0$，有关系
\begin{equation}\tag{11.1.3}
 {}_n\Phi_0(k)\xrightarrow{\sim}{}_nA(K)/{}_nA(K)^0,
 \qquad (n,p)=1.
\end{equation}
而由 ${}_nA_K$ 的\Term{fixed-part}{固定部分}的定义 (2.2.3)，有
\begin{equation}\tag{11.1.4}
 {}_nA(K)^0\xrightarrow{\sim}({}_nA_K)^f(\overline K)
 \xrightarrow{\sim}T_n(A_K)^f(\overline K)\otimes_{\mathbb Z_n}\mathbb Z/n\mathbb Z.
\end{equation}
令 $n$ 为 $\ell$ 的幂，并像 (2.5.3) 那样置
\begin{equation}\tag{11.1.5}
 U=T_\ell(A(\overline K)),\qquad
 \text{从而 }T_\ell(A(\overline K))^f=U^I,
\end{equation}
其中 $U^I$ 是 $U$ 中在\Term{inertia-group}{惯性群} $I$ 下\Term{invariant}{不变}的\Term{submodule}{子模}。于是 (11.1.3) 可以写成
\begin{equation}\tag{11.1.6}
 {}_n\Phi_0(k)\xrightarrow{\sim}(U_n)^I/(U^I)_n,
 \qquad n=\ell^{r+1},
\end{equation}
其中下标 $n$ 表示模 $n$ 的\Term{reduced}{约化}，即在 $\mathbb Z_\ell$ 上与 $\mathbb Z_\ell/n\mathbb Z_\ell$ 作\Term{tensor-product}{张量积}。式 (11.1.6) 两端各群间的过渡\Term{morphism}{态射}都是显然的，因此对 $\ell$ 的幂取\Term{inductive-limit}{归纳极限}，便从 (11.1.6) 得到：

\SourcePage{460}
\begin{SGAresult}{命题}{11.2}
假定 $S$ \Term{strictly-local}{严格局部}，$\ell$ 为不同于 $p$ 的\Term{prime-number}{素数}。则有\Term{canonical}{典范}\Term{isomorphism}{同构}
\begin{equation}\tag{11.2.1}
 \Phi_0(\ell)\xrightarrow{\sim}(U\otimes D_\ell)^I/(U^I\otimes D_\ell),
\end{equation}
其中 $U$ 由 (11.1.5) 定义，且 $D_\ell=\mathbb Q_\ell/\mathbb Z_\ell$。
\end{SGAresult}
因此，只要知道\Term{tate-module}{Tate 模} $U=T_\ell(A(\overline K))$ 在\Term{inertia-group}{惯性群} $I$ \Term{action}{作用}下的 Galois \Term{module-structure}{模结构}，就能明确描述群 $\Phi_0(\ell)$。

\textbf{11.3。}我们将利用显式表达式 (11.2.1)，概述一个\Term{canonical}{典范}\Term{pairing}{配对}
\begin{equation}\tag{11.3.1}
 \Phi_0(\ell)\otimes\Phi'_0(\ell)\longrightarrow\mathbb Q/\mathbb Z
\end{equation}
的定义，它是一个\Term{perfect-duality}{完美对偶}。这里 $\Phi'_0$ 如第1节那样，表示 $A_K$ 的\Term{dual}{对偶} $A'_K$ 的\Term{neron-model}{Néron 模型}所对应的\Term{connected-components}{连通分支群}。几乎可以肯定，这个\Term{pairing}{配对}与 (1.2.1) 定义的\Term{pairing}{配对}至多相差一个\Term{signature}{符号}；不过执笔者没有进行这项核查，读者不妨把证明所需相容性（并找出正确的\Term{signature}{符号}！）作为练习。

置
\[
 U'=T_\ell(A'(\overline K)),\qquad
 T=T_\ell(1)(\overline K)=T_\ell(\overline K^*).
\]
我们有一个与 $I$ 的\Term{action}{作用}相容的、\Term{finite-type}{有限生成}\Term{free-module}{自由} $\mathbb Z_\ell$-模之间的\Term{perfect-pairing}{完美配对}
\begin{equation}\tag{11.3.2}
 \varphi:U\otimes U'\longrightarrow T,
 \qquad (T\simeq\mathbb Z_\ell\text{，非典范地}),
\end{equation}
\SourcePage{461}
并且由 $I$ 的已知结构（第 I 讲\SGAsourcegap{此处第 I 讲的引注编号为空白}），有一个\Term{exact-sequence}{正合列}
\begin{equation}\tag{11.3.3}
 1\longrightarrow R\longrightarrow I\longrightarrow T\longrightarrow1,
\end{equation}
其中 $R$ 是一个阶与 $\ell$ 互素的\Term{profinite-group}{投有限群}。现从数据 (11.3.2)、(11.3.3) 出发，引入群
\begin{equation}\tag{11.3.4}
 \begin{gathered}
 \Phi=(U\otimes D_\ell)^I/(U^I\otimes D_\ell),\qquad
 \Phi'=(U'\otimes D_\ell)^I/(U'^I\otimes D_\ell),\\
 D_\ell=\mathbb Q_\ell/\mathbb Z_\ell.
 \end{gathered}
\end{equation}
我们将定义一个\Term{perfect-duality}{完美对偶}
\begin{equation}\tag{11.3.5}
 \Phi\otimes\Phi'\longrightarrow\mathbb Q/\mathbb Z,
\end{equation}
特别地，它将给出所宣称的\Term{pairing}{配对} (11.3.1)。考虑 $I$-模的\Term{exact-sequence}{正合列}
\begin{equation}\tag{11.3.6}
 0\longrightarrow U\xrightarrow{i}U\otimes_{\mathbb Z_\ell}\mathbb Q_\ell
 \xrightarrow{j}U\otimes D_\ell\longrightarrow0,
\end{equation}
及其相应的\Term{cohomology}{上同调}\Term{exact-sequence}{正合列}；这里 $I$ 的\Term{cohomology}{上同调}由取值于所考察\Term{topological-group}{拓扑群}中的\Term{continuous-cochain}{连续上链}定义：
\begin{equation}\tag{11.3.7}
 \begin{aligned}
 0&\longrightarrow H^0(I,U)\xrightarrow{i^0}H^0(I,U\otimes\mathbb Q_\ell)
 \xrightarrow{j^0}H^0(I,U\otimes D_\ell)\\
 &\xrightarrow{\partial}H^1(I,U)\xrightarrow{i^1}H^1(I,U\otimes\mathbb Q_\ell).
 \end{aligned}
\end{equation}
其中相应的自然认同为
\[
 \begin{aligned}
 H^0(I,U)&=U^I,\\
 H^0(I,U\otimes\mathbb Q_\ell)&=(U\otimes\mathbb Q_\ell)^I
 \simeq U^I\otimes_{\mathbb Z}\mathbb Q,\\
 H^0(I,U\otimes D_\ell)&=(U\otimes D_\ell)^I,\\
 H^1(I,U\otimes\mathbb Q_\ell)&\simeq H^1(I,U)\otimes_{\mathbb Z}\mathbb Q.
 \end{aligned}
\]
显然，(11.3.4) 定义的 $\Phi$ 正是 $\operatorname{coker}j^0$，故\Term{canonical}{典范}\Term{isomorphism}{同构}于 $\operatorname{Ker}i^1$；即有\Term{canonical}{典范}\Term{exact-sequence}{正合列}
\begin{equation}\tag{11.3.8}
 0\longrightarrow\Phi\longrightarrow H^1(I,U)
 \longrightarrow H^1(I,U)\otimes_{\mathbb Z}\mathbb Q,
\end{equation}
\SourcePage{462}
这表明，$\Phi$ \Term{canonical}{典范}\Term{isomorphism}{同构}于 $H^1(I,U)$ 的\Term{torsion-subgroup}{挠子群}。

另一方面，由\Term{exact-sequence}{正合列} (11.3.3)，对每个 $n=\ell^r$ 都得到\Term{isomorphism}{同构}
\begin{equation}\tag{11.3.9}
 H^1(I,T_n)\xrightarrow{\sim}\operatorname{Hom}(I,T_n)
 \xrightarrow{\sim}\operatorname{Hom}(T,T_n)\xrightarrow{\sim}\mathbb Z/n\mathbb Z.
\end{equation}
（注意，我们让 $I$ \Term{trivial-action}{平凡地作用}于 $T$，从而也\Term{trivial-action}{平凡地作用}于 $T_n=T/nT$。）因此，对于每个\Term{finite}{有限} $I$-$\mathbb Z_\ell$-模 $M$，若
\begin{equation}\tag{11.3.10}
 M'=\operatorname{Hom}(M,T\otimes D_\ell)
\end{equation}
表示它的“\Term{cartier-dual}{Cartier 对偶}”，就有\Term{canonical}{典范}\Term{pairing}{配对}
\begin{equation}\tag{11.3.11}
 H^i(I,M)\otimes H^{1-i}(I,M')\longrightarrow H^1(I,T\otimes D_\ell)
 \xrightarrow{\sim}D_\ell,
\end{equation}
其中最后的\Term{isomorphism}{同构}由 (11.3.9) 中的\Term{isomorphism}{同构}对 $\ell$ 的幂 $n$ 取\Term{inductive-limit}{归纳极限}得到。这些\Term{pairing}{配对}是\Term{perfect-duality}{完美对偶}，这是熟知的（例如参见 SGA 5，第 I 讲5）。在这些\Term{pairing}{配对}中，取 $M=U_n$，从而 $M'_n=U'_n$，取\Term{limit}{极限}便得到一个\Term{perfect-pairing}{完美配对}
\begin{equation}\tag{11.3.12}
 H^0(I,U\otimes D_\ell)\times H^1(I,U')\longrightarrow D_\ell,
\end{equation}
其中左端的第二个因子是\Term{finite-type}{有限生成} $\mathbb Z_\ell$-模。对\Term{pairing}{配对} (11.3.12) 再取一次\Term{limit}{极限}（过渡\Term{morphism}{态射}形如 $\ell^r\operatorname{id}$），便得到\Term{finite-dimensional}{有限维} $\mathbb Q_\ell$-\Term{vector-space}{向量空间}之间的第二个\Term{perfect-pairing}{完美配对}
\begin{equation}\tag{11.3.13}
 H^0(I,U\otimes\mathbb Q_\ell)\times H^1(I,U'\otimes\mathbb Q_\ell)
 \longrightarrow D_\ell.
\end{equation}
以 $i',j'$ 表示将 (11.3.6) 中 $U$ 换为 $U'$ 后定义的\Term{morphism}{态射}，
\SourcePage{463}
便容易看出，以下两个\Term{homomorphism}{同态}
\[
 H^0(I,U\otimes\mathbb Q_\ell)\xrightarrow{j^0}H^0(I,U\otimes D_\ell),
\]
\[
 H^1(I,U'\otimes\mathbb Q_\ell)\xleftarrow{i'^1}H^1(I,U')
\]
相对于\Term{pairing}{配对} (11.3.12)、(11.3.13) 互为\Term{transpose}{转置}。立即可知，$\operatorname{coker}j^0$ 与 $\operatorname{Ker}i'^1$ 通过一个取值于 $D_\ell$ 的\Term{canonical}{典范}\Term{pairing}{配对}互为\Term{dual}{对偶}；这就是所宣称的\Term{pairing}{配对} (11.3.5)。

\textbf{11.4。}上面借助\Term{connected-components}{连通分支群}的 $\ell$-\Term{primary-component}{准素分量}的表达式 (11.2.1)，对\Term{pairing}{配对} (11.3.1) 所作的描述，只适用于 $\ell\ne p$。现假定 $A_K$ 在 $S$ 上具有\Term{semistable-reduction}{半稳定约化}，我们还将定义\Term{pairing}{配对}
\begin{equation}\tag{11.4.1}
 \Phi_0(\ell)\otimes\Phi'_0(\ell)\longrightarrow\mathbb Q/\mathbb Z,
\end{equation}
它们是\Term{perfect-duality}{完美对偶}，且对 $\ell$ 没\Term{finite}{有限}制。像11.4中那样，还应证明这里重新得到的是 (1.2.1) 定义的\Term{pairing}{配对}（也许相差一个有待确定的\Term{signature}{符号}）；如此便可在\Term{semistable-reduction}{半稳定约化}情形，以及限于 $\ell\ne p$ 的 $\ell$-\Term{primary-component}{准素分量}的情形，证明猜想1.3。

为定义 (11.4.1)，我们本质上将给出一个用 $K$ 上的\Term{barsotti-tate-group}{Barsotti--Tate 群} $T_\ell(A_K)$ 表示 $\Phi_0(\ell)$ 的表达式，
\SourcePage{464}
它在当前情形将发挥11.2在先前情形中的同样\Term{action}{作用}：

\begin{SGAresult}{\Term{theorem}{定理}}{11.5}
假定 $S$ 是\Term{henselian}{Hensel}的，且 $A_K$ 在 $S$ 上具有\Term{semistable-reduction}{半稳定约化}。设 $\ell$ 为\Term{prime-number}{素数}，考虑\Term{monodromy-pairing}{单值性配对} (9.1.2)，以及由此得到的\Term{canonical}{典范}\Term{homomorphism}{同态}
\[
 \widetilde u_\ell:\underline M_\ell\longrightarrow(\underline M'_\ell)^\vee.
\]
置 $\Phi(\ell)=\operatorname{coker}\widetilde u_\ell$；则有\Term{canonical}{典范}\Term{isomorphism}{同构}
\begin{equation}\tag{11.5.1}
 \Phi_0(\ell)\xrightarrow{\sim}\Phi(\ell)\times_Ss.
\end{equation}
\end{SGAresult}
这个陈述立刻给出\Term{perfect-pairing}{完美配对} (11.4.1) 的定义，因为对 $A'_K$ 而非 $A_K$ 应用11.5，并利用 (10.2.2)，还能得到\Term{canonical}{典范}\Term{isomorphism}{同构}
\[
 \Phi'_0(\ell)\xrightarrow{\sim}\Phi'(\ell)\times_Ss,
 \qquad \Phi'(\ell)=\operatorname{coker}
 \bigl({}^t\widetilde u_\ell:\underline M'_\ell\longrightarrow\underline M_\ell^\vee\bigr).
\]

\textbf{注记 11.5.2。}a) \Term{theorem}{定理}11.5特别蕴含，$\widetilde u_\ell$ 是\Term{isogeny}{同源}，即 $u_\ell$ 是\Term{separating-pairing}{分离的}；这也包括 $\ell=p$ 的情形，在此情形，该断言并非由9.6的定义显然可得。注意，这里不使用将于下一节证明的、深得多的\Term{theorem}{定理}10.4。

b) 利用10.4所给出的\Term{monodromy-pairing}{单值性配对} $u:\underline M\otimes\underline M'\longrightarrow\mathbb Z_S$，并以 $\Phi$ 表示相应\Term{homomorphism}{同态}
\[
 \widetilde u:\underline M\longrightarrow(\underline M')^\vee
\]
的\Term{cokernel}{余核}，
\SourcePage{465}
便可用\Term{canonical}{典范}\Term{isomorphism}{同构}
\begin{equation}\tag{11.5.3}
 \Phi_0\xrightarrow{\sim}\Phi\times_Ss
\end{equation}
得到11.5的一个不依赖于 $\ell$ 的\Term{equivalence}{等价}表述。

\textbf{11.6。}为定义\Term{isomorphism}{同构} (11.5.1)，可限于 $S$ \Term{strictly-local}{严格局部}的情形（11.0）；于是 $\underline M,\underline M'$ 都是\Term{constant-group}{常值}的，分别取值于
\begin{equation}\tag{11.6.1}
 M=D(T'_0),\qquad M'=D(T_0),
\end{equation}
其中 $T_0,T'_0$ 分别是 $A_0,A'_0$ 的最大\Term{torus}{环面}。同样将 $\Phi_0$ 认同为通常意义下的\Term{finite-group}{有限群}，便须在它的 $\ell$-\Term{primary-component}{准素分量} $\Phi_0(\ell)$ 与
\[
 \Phi(\ell)=\operatorname{coker}\bigl(\widetilde u_\ell:M_\ell\longrightarrow(M'_\ell)^\vee\bigr)
\]
之间定义\Term{canonical}{典范}\Term{isomorphism}{同构}。事实上，我们将直接定义一个\Term{canonical}{典范}\Term{isomorphism}{同构}
\begin{equation}\tag{11.6.2}
 \Phi_0(\ell)\xrightarrow{\sim}
 \ker\bigl(\widetilde u_\ell\otimes D_\ell:
 M_\ell\otimes D_\ell\longrightarrow(M'_\ell)^\vee\otimes D_\ell\bigr).
\end{equation}
由于 $\Phi_0(\ell)$ 是\Term{finite}{有限}的，且 $M_\ell,(M'_\ell)^\vee$ 的\Term{rank}{秩}相同，首先可知 $\widetilde u_\ell:M_\ell\longrightarrow(M'_\ell)^\vee$ 是\Term{isogeny}{同源}；进而可知 (11.6.2) 的右端\Term{canonical}{典范}\Term{isomorphism}{同构}于 $\Phi(\ell)$：
\begin{equation}\tag{11.6.3}
 \Phi(\ell)=\operatorname{coker}\widetilde u_\ell
 \xrightarrow{\sim}\ker(\widetilde u_\ell\otimes D_\ell).
\end{equation}
这是因为，通过 $\widetilde u_\ell$ 将 $M_\ell$ 认同为 $(M'_\ell)^\vee$ 的\Term{finite-index-submodule}{有限指数子模}后，可以写成
\[
 \Phi(\ell)=V/M_\ell,\qquad \Phi'(\ell)=V/(M'_\ell)^\vee,
\]
其中 $V=M_\ell\otimes\mathbb Q_\ell\simeq(M'_\ell)^\vee\otimes\mathbb Q_\ell$。将 (11.6.2)、(11.6.3) 复合，立刻得到待定义的\Term{isomorphism}{同构} (11.5.1)，故一切归结为定义\Term{canonical}{典范}\Term{isomorphism}{同构} (11.6.2)。

\SourcePage{466}
将 $M_\ell\otimes D_\ell$（相应地，$(M'_\ell)^\vee\otimes D_\ell$）写成各 $M_n$（相应地，$(M'_n)^\vee$）的\Term{inductive-limit}{归纳极限}，其中 $n$ 遍历 $\ell$ 的幂，可知定义 (11.6.2) \Term{equivalence}{等价}于定义\Term{canonical}{典范}\Term{isomorphism}{同构}
\begin{equation}\tag{11.6.4}
 {}_n\Phi_0\xrightarrow{\sim}\operatorname{Ker}
 \bigl(\widetilde u_n:M_n\longrightarrow(M'_n)^\vee\bigr),
 \qquad n=\ell^{r+1},
\end{equation}
并要求它们与 $n$ 变化时显然的过渡\Term{morphism}{态射}相容。（这里稍微滥用记号，$\widetilde u_n$ 现在表示从 $\widetilde u_\ell:M_\ell\longrightarrow(M'_\ell)^\vee$ 模 $n$ \Term{reduced}{约化}得到的\Term{homomorphism}{同态}。）

为此，考虑 $S$ 上的\Term{group-scheme}{群概形} ${}_nA$。由\Term{semistable-reduction}{半稳定约化}假设，它在 $S$ 上\Term{flat}{平坦}且\Term{quasi-finite}{拟有限}（2.2.1）。有一个\Term{canonical}{典范}\Term{filtration}{滤过}
\begin{equation}\tag{11.6.5}
 {}_nA\supset({}_nA^0)^f\supset({}_nA^0)^t\supset0,
\end{equation}
其中上标 $f$（相应地，上标 $t$）表示所标记\Term{group-scheme}{群概形}的“\Term{fixed-part}{固定部分}”（2.2.3）（相应地，“\Term{toric-part}{环面部分}”（5.1））。对于5.1中\Term{toric-part}{环面部分}的定义，应注意当时假定了 $S$ \Term{complete}{完备}，而这里也可如此假定（11.0）。（此外，6.1已经表明，如果愿意，此构造可以推广到 $S$ 为\Term{henselian}{Hensel}的情形……）置
\begin{equation}\tag{11.6.6}
 {}_n\Psi={}_nA/({}_nA^0)^f;
\end{equation}
这是 $S$ 上一个\Term{flat}{平坦}、\Term{separated}{分离}且\Term{quasi-finite}{拟有限}的\Term{group-scheme}{群概形}（SGA 3，第 VI$_A$ 讲3.2及第 VI$_B$ 讲9.2），其\Term{generic-fiber}{泛纤维}为
\[
 {}_n\Psi_K=({}_nA_K)/({}_nA_K)^f,
\]
\SourcePage{467}
其中，按定义，$({}_nA_K)^f=({}_nA^0)^f_K$ 表示 Barsotti--Tate 的\Term{pro-group}{pro-群} $T_\ell(A_K)^f$ 的第 $r$ 个\Term{set-component}{分量}（$n=\ell^{r+1}$）。关系 (5.5.8) 的第一个给出\Term{canonical}{典范}\Term{isomorphism}{同构}
\begin{equation}\tag{11.6.7}
 {}_n\Psi_K\xrightarrow{\sim}\underline M_{nK},
 \qquad \underline M_n\mathrel{\mathop=^{\mathrm{def}}}
 \underline M\otimes_{\mathbb Z}(\mathbb Z/n\mathbb Z).
\end{equation}
事实上，有如下结果：
\begin{SGAresult}{\Term{lemma}{引理}}{11.6.8}
$S$ 上的\Term{group-scheme}{群概形} ${}_n\Psi$ (11.6.6) 在 $S$ 上\Term{etale}{平展}，而\Term{isomorphism}{同构} (11.6.7) 唯一地延拓为一个\Term{morphism}{态射}
\begin{equation}\tag{11.6.8.1}
 {}_n\Psi\lhook\joinrel\longrightarrow\underline M_n,
\end{equation}
它是一个\Term{open-immersion}{开浸入}。
\end{SGAresult}
\begin{proof}
由 (11.6.7)，要说明 ${}_n\Psi$ \Term{etale}{平展}，只须核实它在\Term{special-fiber}{特殊纤维}的各点处\Term{etale}{平展}，也就是核实该\Term{special-fiber}{特殊纤维}\Term{etale}{平展}。这显然成立，因为该\Term{fiber}{纤维}为 ${}_nA_0/{}_nA_0^0\subset A_0/A_0^0=\Phi_0$。因此 ${}_n\Psi$ 是\Term{normal}{正规}\Term{scheme}{概形}；又因 $\underline M_n$ 在 $S$ 上\Term{finite}{有限}，延拓 (11.6.8.1) 的存在性立即得到。这是 $S$ 上的\Term{etale}{平展}\Term{group-scheme}{群概形}之间的\Term{morphism}{态射}，故其\Term{kernel}{核}是 $S$ 上的一个\Term{etale}{平展}群。这个\Term{kernel}{核}像 ${}_n\Psi$ 一样\Term{separated}{分离}，且由 (11.6.7)，它的\Term{generic-fiber}{泛纤维}为\Term{unit-group}{单位群}，故它本身为\Term{unit-group}{单位群}（3.1.3）。因此 (11.6.8.1) 是一个\Term{etale}{平展}\Term{monomorphism}{单态射}，从而是\Term{open-immersion}{开浸入}（EGA IV 17.9.1）。\Term{lemma}{引理}得证。
\end{proof}
现在置
\begin{equation}\tag{11.6.9}
 {}_n\Phi=({}_nA)^f/({}_nA^0)^f\subset{}_n\Psi.
\end{equation}
于是 ${}_n\Phi$ 是 ${}_n\Psi$ 的一个\Term{finite}{有限}\Term{open-subscheme}{开子概形}，因为 $({}_nA)^f$ 是 ${}_nA$ 的\Term{finite}{有限}\Term{open-subgroup}{开子群}\Term{scheme}{概形}；
\SourcePage{468}
更确切地说，立即可见
\begin{equation}\tag{11.6.10}
 {}_n\Phi\xrightarrow{\sim}({}_n\Psi)^f.
\end{equation}
复合 (11.6.9) 与 (11.6.8.1)，得到一个\Term{open-immersion}{开浸入}
\begin{equation}\tag{11.6.11}
 {}_n\Phi\lhook\joinrel\longrightarrow\underline M_n.
\end{equation}
另一方面，有
\[
 {}_n\Phi\times_Ss={}_n\Psi\times_Ss={}_nA_0/{}_nA_0^0.
\]
由于 $A_0^0$ 是 $n$-\Term{divisible}{可除}的（2.2.1），立刻可知最后一项可认同于 ${}_n\Phi_0$；故得到一个\Term{isomorphism}{同构}
\begin{equation}\tag{11.6.12}
 {}_n\Phi\times_Ss\xrightarrow{\sim}{}_n\Phi_0,
\end{equation}
它表明，${}_n\Phi$ 被定义为\Term{special-fiber}{特殊纤维}为 ${}_n\Phi_0$ 的、$S$ 上的\Term{finite}{有限}\Term{etale}{平展}\Term{scheme}{概形}，并且在唯一\Term{isomorphism}{同构}的意义下确定。因此，为通过\Term{open-immersion}{开浸入} (11.6.11) 定义所求\Term{isomorphism}{同构} (11.6.4)，只须证明：
\begin{SGAresult}{\Term{lemma}{引理}}{11.6.13}
\Term{canonical}{典范}\Term{homomorphism}{同态} (11.6.11) 的像等于 $\widetilde u_n:\underline M_n\longrightarrow(\underline M'_n)^\vee$ 的\Term{kernel}{核}。
\end{SGAresult}
此外，所需的各\Term{homomorphism}{同态} (11.6.11) 与 $n$ 变化时过渡\Term{homomorphism}{同态}之间的相容性，由定义显然成立。所以有了11.6.13，11.5的证明就完成了。

\SourcePage{469}
\begin{proof}[\Term{lemma}{引理}11.6.13的证明]
注意，$K_n=\ker\widetilde u_n$ 也是\Term{canonical}{典范}\Term{pairing}{配对}
\[
 u_n:\underline M_n\times\underline M'_n\longrightarrow(\mathbb Z/n\mathbb Z)_S
\]
的左\Term{annihilator}{零化子}。它的显式定义已在9.4与9.6中给出，所依据的是\Term{extension}{扩张}
\begin{equation}\tag{11.6.14}
 \begin{gathered}
 0\longrightarrow(\underline M'_n)^\vee(1)\longrightarrow F_n
 \longrightarrow R_n\longrightarrow0,\qquad F_n=({}_nA^0)^f,\\
 0\longrightarrow\underline M_n^\vee(1)\longrightarrow F'_n
 \longrightarrow R'_n\longrightarrow0,\qquad F'_n=({}_nA'^0)^f
 \end{gathered}
\end{equation}
（其中 $R_n,R'_n$ 互为\Term{cartier-dual}{Cartier 对偶}），以及 $E_{n\eta}$ 被 $F_{n\eta}$ 的\Term{blended-extension}{混合扩张} $X_{n\eta}$，后者是 (11.7.5) 的\Term{generic-fiber}{泛纤维}。这里置
\begin{equation}\tag{11.6.15}
 E_n=D(F_n),\qquad
 0\longrightarrow\underbrace{R_n}_{\simeq D(R'_n)}\longrightarrow E_n
 \longrightarrow\underline M_n\longrightarrow0.
\end{equation}
由9.4.6立刻可知，$K_n$ 是 $\underline M_n$ 中满足下述条件的最大\Term{finite}{有限}\Term{etale}{平展}\Term{subgroup}{子群} $K'_n$：$K'_n$ 在\Term{blended-extension}{混合扩张} $X_{n\eta}$ 中的\Term{inverse-image}{逆像} $\overline X_{n\eta}$，视为 $\overline E_{n\eta}$ 被 $F_{n\eta}$ 的\Term{blended-extension}{混合扩张}时，可以延拓为 $\overline E_n$ 被 $F_n$ 的\Term{blended-extension}{混合扩张}；这里 $\overline E_n$ 是所指的那个\Term{subgroup}{子群}，它是 $K_n$ 被 $R_n$ 的\Term{extension}{扩张}。而 $\underline M_n$ 的\Term{subgroup}{子群} ${}_n\Phi$ (11.6.11) 显然满足此条件，因为在 $S$ 上已有相应的\Term{blended-extension}{混合扩张} $({}_nA)^f$。因此 ${}_n\Phi\subset K_n$。为证明相反的包含关系，考虑上述构造中取 $K'_n=K_n$ 时定义的\Term{blended-extension}{混合扩张} $\overline X_{n\eta}$ 的延拓
\SourcePage{470}
$\overline X_n$，它是 $\overline E_n$ 被 $F_n$ 的\Term{blended-extension}{混合扩张}。于是它是一个\Term{extension}{扩张}
\[
 0\longrightarrow F_n\longrightarrow\overline X_n\longrightarrow K_n\longrightarrow0,
\]
且有\Term{canonical}{典范}\Term{homomorphism}{同态}
\begin{equation}\tag{*}
 F_n=({}_nA^0)^f\lhook\joinrel\longrightarrow A,
\end{equation}
其\Term{generic-fiber}{泛纤维}由\Term{canonical}{典范}\Term{homomorphism}{同态}
\begin{equation}\tag{**}
 \overline X_{n\eta}\lhook\joinrel\longrightarrow X_{n\eta}
 ={}_nA_K\lhook\joinrel\longrightarrow A_K
\end{equation}
诱导。由 (5.9.2)，立即可知\Term{homomorphism}{同态} (**) 可唯一延拓为一个\Term{homomorphism}{同态}
\begin{equation}\tag{***}
 \overline X_n\longrightarrow A,
\end{equation}
它诱导 (*)。由于 (**) 将\Term{generic-fiber}{泛纤维}映入 ${}_nA$，而后者是 $A$ 的\Term{closed-subscheme}{闭子概形}，便知该\Term{homomorphism}{同态}经由 $\overline X_n\longrightarrow{}_nA$ \Term{factorization}{分解}；又因 $\overline X_n$ \Term{finite}{有限}，它经由 $({}_nA)^f$ \Term{factorization}{分解}。考察在\Term{generic-fiber}{泛纤维}上诱导的\Term{homomorphism}{同态}，得到
\[
 \overline X_{n\eta}\subset({}_nA)^f_\eta,
 \quad\text{即 }K_{n\eta}\subset{}_n\Phi_\eta,
 \quad\text{即 }K_n\subset{}_n\Phi.
\]
证毕。
\end{proof}

\textbf{注记 11.7。}11.7中的构造表明，知道 $S$ 上的\Term{extension}{扩张} (11.6.14)，以及 $E_n$ 被 $F_n$ 的\Term{blended-extension}{混合扩张} $X_{n\eta}=({}_nA_K)$，就可以在唯一\Term{isomorphism}{同构}的意义下重建\Term{blended-extension}{混合扩张} (11.6.5)，从而也重建 ${}_nA$ 的\Term{subgroup}{子群} $({}_nA)^f\supset({}_nA^0)^f$。事实上，由9.4可得\Term{pairing}{配对} $u_n:\underline M_n\times\underline M'_n\longrightarrow(\mathbb Z/n\mathbb Z)_S$，而上面的构造
\SourcePage{471}
给出 $({}_nA)^f$，将其表为 $\overline E_n$（即 $\operatorname{Ker}\widetilde u_n$ 经 $E_n\longrightarrow\underline M_n$ 的\Term{inverse-image}{逆像}）被 $F_n$ 的\Term{blended-extension}{混合扩张} $\overline X_n$；它延拓由 $X_{n\eta}$ 诱导的\Term{blended-extension}{混合扩张} $\overline X_{n\eta}$。注意，由 (9.4.8)，\Term{blended-extension}{混合扩张} $\overline X_n$ 由此在唯一\Term{isomorphism}{同构}的意义下确定。然后，将 $\overline X_n$ 与 $X_{n\eta}$ 沿共同\Term{open-subset}{开集} $\overline X_{n\eta}$ \Term{gluing}{粘合}，便得到 ${}_nA$。

因此，若知道结构要素 (9.6.1)，以及 $E_\eta$ 被 $F_\eta$ 的\Term{blended-extension}{混合扩张} $X_\eta=T_\ell(A_K)$（这\Term{equivalence}{等价}于对每个形如 $\ell^{r+1}$、$r\ge0$ 的 $n$，知道\Term{extension}{扩张} (11.6.14) 及 $X_{n\eta}$），就能重建\Term{blended-extension}{混合扩张} (11.6.5) 的\Term{inductive-system}{归纳系}（或者\Term{projective-system}{投射系}，可任选其一）。

\textbf{11.8。}像8.3中那样，将 $K$ 的\Term{separable-closure}{可分闭包} $\overline K$ 视为其\Term{finite}{有限}子\Term{extension}{扩张} $K'$ 的\Term{inductive-limit}{归纳极限}。对每个这样的 $K'$，考虑 $A_{K'}=A_K\otimes_KK'$ 的\Term{neron-model}{Néron 模型}
\[
 A[K'],
\]
由此得到一个\Term{finite}{有限}\Term{etale}{平展}的\Term{connected-components}{连通分支群}
\begin{equation}\tag{11.8.1}
 \Phi_0[K']=A[K']_0/A[K']_0^0,
\end{equation}
它定义在 $S$ 在 $K'$ 中的\Term{normalization}{正规化} $S[K']$ 的\Term{residue-field}{剩余域}上。为简便起见，假定 $S$ 的\Term{residue-field}{剩余域}\Term{separably-closed}{可分闭}，从而各 $S[K']$ 的\Term{residue-field}{剩余域}也如此。于是各群 (11.8.1) 可认同为通常意义下的\Term{finite-group}{有限群} $\Phi[K']$；显然，当 $K'$ 变化时，它们构成一个\Term{inductive-system}{归纳系}，我们正要研究这个\Term{inductive-system}{归纳系}。我们特别关心群
\SourcePage{472}
\begin{equation}\tag{11.8.2}
 \Psi=\varinjlim_{K'}\Phi[K'].
\end{equation}
由 (3.6) 已知，对充分大的 $K'$，譬如 $K'\supset K'_0$，$A[K']$ 具有\Term{semistable-reduction}{半稳定约化}；因此当 $K'\supset K'_0$ 时，过渡\Term{morphism}{态射}是\Term{injective}{单射}（3.3）。另外，采用5.5的记号并沿用于此，对 $K'\supset K'_0$ 有\Term{canonical}{典范}\Term{isomorphism}{同构}
\begin{equation}\tag{11.8.3}
 \Phi[K'](\ell)\xrightarrow{\sim}
 \ker\bigl(\widetilde u[K']_\ell\otimes D_\ell:
 M_\ell\otimes D_\ell\longrightarrow(M'_\ell)^\vee\otimes D_\ell\bigr),
\end{equation}
这是对 $A[K']$ 应用 (11.6.2) 所得。又可立刻核实，相应的\Term{immersion}{浸入}
\[
 \Phi[K'](\ell)\lhook\joinrel\longrightarrow M_\ell\otimes D_\ell
\]
与 $K'$ 变化时的过渡\Term{morphism}{态射}相容，因此取\Term{limit}{极限}就得到 $\Psi$ 的 $\ell$-\Term{primary-component}{准素分量}的一个\Term{canonical}{典范}\Term{monomorphism}{单态射}
\begin{equation}\tag{11.8.4}
 \Psi(\ell)\mathrel{\mathop=^{\mathrm{def}}}
 \varinjlim_{K'}\Phi[K'](\ell)\lhook\joinrel\longrightarrow M_\ell\otimes D_\ell.
\end{equation}
对所有 $\ell$ 求和，得到
\begin{equation}\tag{11.8.5}
 \Psi\mathrel{\mathop=^{\mathrm{def}}}\varinjlim_{K'}\Phi[K']
 \longrightarrow M\otimes\mathbb Q/\mathbb Z.
\end{equation}

\begin{SGAresult}{\Term{theorem}{定理}}{11.9}
\Term{canonical}{典范}\Term{homomorphism}{同态} (11.8.5) 是\Term{isomorphism}{同构}。
\end{SGAresult}
\begin{proof}
已经知道它是\Term{monomorphism}{单态射}；还须证明它是\Term{isomorphism}{同构}。为此，要利用 (11.8.3) 对 $\Phi[K'](\ell)$ 的像的刻画，以及 $K'$ 变化时 $\widetilde u(K')$ 的变化规律。此规律由 (10.3.4) 给出。
\SourcePage{473}
归结到 $S$ \Term{strictly-local}{严格局部}的情形是允许的；于是，证明11.9就归结为下述熟知陈述（将它应用于 $S(K'_0)$）：若 $S$ 是\Term{fraction-field}{分式域}为 $K$ 的\Term{strictly-local}{严格局部}\Term{trait}{离散赋值环谱}，则对于每个整数 $n\ge1$，存在 $K$ 的一个\Term{separable-extension}{可分扩张} $K'$，其\Term{degree}{次数}是 $n$ 的倍数（参见第 I 讲\SGAsourcegap{此处第 I 讲的引注编号为空白}）。
\end{proof}

\begin{SGAresult}{推论}{11.10}
在8.3的条件下（$S$ \Term{complete}{完备}），并采用该处记号，$K$ 上的\Term{ind-finite}{ind-有限}\Term{group-scheme}{群概形} $A_K^\flat/A_K^{\flat0}$，即\Term{group-scheme}{群概形} $A_K^\flat$ 的\Term{connected-components}{连通分支群}\Term{scheme}{概形}，\Term{canonical}{典范}\Term{isomorphism}{同构}于 $\underline M_K\otimes\mathbb Q/\mathbb Z$，其中 $\underline M_K$ 是 (8.1.8) 中定义的 $K$ 上的\Term{twisted-constant-group}{扭常值群}。
\end{SGAresult}
\begin{proof}
考虑到 $\underline M_K$ 的定义，这是 (8.2.4) 与11.9的直接推论。
\end{proof}

\textbf{注记 11.11。}a) 当仅考虑 $\ell\ne p$ 的 $\ell$-\Term{primary-component}{准素分量}时，\Term{theorem}{定理}11.9由11.2立刻可得，无须诉诸11.6中那些较为精细的论述（经由9.4）。同样，当假定 $K$ 的\Term{characteristic}{特征}为零时，证明11.9也无须使用 (11.8.3)。这是因为，${}_n\Phi(K')(\ell)$ 在 $M_n={}_n(M_\ell\otimes D_\ell)$ 中的像显然包含 ${}_nA(K')$ 的像；后者是 ${}_nA$ 的 $K'$-值点群，当 $K'$ 充分大时，它等于 ${}_nA(\overline K)$，因而\Term{surjective}{满射}到 $M_n$。当 $\operatorname{car}K=p$ 时，对 $p$-\Term{primary-component}{准素分量}的11.10似乎不存在“平凡的”（即不使用9.4的）证明。

b) 存在\Term{isomorphism}{同构} $\Psi\xrightarrow{\sim}(\mathbb Q/\mathbb Z)^r$ 这一事实，以及 $r$ 的正确数值，早在1964年就已由 J.-P.~Serre 猜出。
